\documentclass[11pt,a4paper]{article}
\usepackage[utf8]{inputenc}
\usepackage[T1]{fontenc}
\usepackage[catalan]{babel}
\usepackage{lmodern}
\usepackage[sfdefault,lining]{FiraSans}
\usepackage[a4paper,top=2.3cm,bottom=2.2cm,left=1.8cm,right=1.8cm,headheight=1.3cm,headsep=0.4cm,footskip=1cm]{geometry}
\usepackage{amsmath,amssymb}
\usepackage{icomma}
\usepackage{graphicx}
\usepackage[table]{xcolor}
\usepackage{tikz}
\usepackage{fontawesome5}
\usepackage{ccicons}
\usepackage{fancyhdr}
\usepackage{lastpage}
\usepackage{tasks}
\usepackage{needspace}
\usepackage[most]{tcolorbox}
\usepackage{array}
\usepackage[hidelinks,pdftitle={Fitxa d'exercicis de la unitat 1: Nombres racionals. 3r d'ESO},pdfsubject={Matemàtiques 3r ESO},pdfkeywords={nombres enters, fraccions, nombres racionals, decimals, fracció generatriu}]{hyperref}

\definecolor{navy}{HTML}{1B2748}
\definecolor{accent}{HTML}{E8603A}
\definecolor{teal}{HTML}{1F9E90}
\definecolor{soft}{HTML}{EEF1F8}
\definecolor{muted}{HTML}{5D6685}

\setlength{\parindent}{0pt}
\setlength{\parskip}{3pt}
\settasks{label=\alph*), label-format=\bfseries\color{navy}, label-width=1.3em, item-indent=1.8em, column-sep=0.6em, after-item-skip=2pt, before-skip=3pt, after-skip=4pt}

% ---------- capçalera i peu ----------
\newcommand{\llicencia}{\href{https://creativecommons.org/licenses/by-nc-sa/4.0/deed.ca}{\ccbyncsa}\ \ Llicència Creative Commons Reconeixement-NoComercial-CompartirIgual 4.0 Internacional (CC BY-NC-SA 4.0)}
\pagestyle{fancy}
\fancyhf{}
\fancyhead[L]{\raisebox{-0.35\height}{\includegraphics[height=0.95cm]{logo-gva.png}}}
\fancyhead[C]{\small\color{muted}Fitxa d'exercicis · Unitat 1: Nombres racionals · 3r d'ESO}
\fancyhead[R]{\raisebox{-0.35\height}{\includegraphics[height=0.95cm]{logo-vermellar.png}}}
\fancyfoot[L]{\scriptsize\color{muted}\llicencia}
\fancyfoot[R]{\small\color{muted}\thepage}
\renewcommand{\headrulewidth}{0.4pt}
\renewcommand{\headrule}{\color{muted!40}\hrule height \headrulewidth}
\fancypagestyle{primera}{\fancyhf{}\renewcommand{\headrulewidth}{0pt}
  \fancyfoot[L]{\scriptsize\color{muted}\llicencia}
  \fancyfoot[R]{\small\color{muted}\thepage}}

% ---------- nivells i exercicis ----------
\newcommand{\estrelles}[1]{{\color{accent}\ifcase#1\or\faStar\faStar[regular]\faStar[regular]\or\faStar\faStar\faStar[regular]\or\faStar\faStar\faStar\fi}}
\newcounter{exer}
\newcommand{\ex}[1]{\par\needspace{4\baselineskip}\medskip\stepcounter{exer}\noindent{\color{navy}\bfseries\theexer.}\ \ {\small\estrelles{#1}}\ \ }

\newcommand{\seccio}[2]{\par\needspace{14\baselineskip}\bigskip
  \begin{tcolorbox}[colback=navy,colframe=navy,coltext=white,arc=2mm,boxrule=0pt,left=3mm,right=3mm,top=1.2mm,bottom=1.2mm]
  {\large\bfseries #1\quad #2}
  \end{tcolorbox}\vspace{-2pt}}

% ---------- recta real ----------
\newcommand{\rectabuida}{\begin{tikzpicture}[x=0.42cm,baseline=0]
  \draw[-latex,muted] (-6.7,0)--(6.8,0);
  \foreach \i in {-6,...,6}{\draw[muted] (\i,0.12)--(\i,-0.12) node[below=-1pt,font=\scriptsize,text=muted]{$\i$};}
  \end{tikzpicture}}
% \rectaint{x1}{x2}{estil esquerra}{estil dreta}: estil = ple, buit, inf
\newcommand{\extrem}[2]{\ifx#2p\fill[accent] (#1,0) circle (0.13);\else\ifx#2b\filldraw[fill=white,draw=accent,line width=1.2pt] (#1,0) circle (0.13);\fi\fi}
\newcommand{\rectaint}[4]{\begin{tikzpicture}[x=0.42cm,baseline=0]
  \draw[-latex,muted] (-6.7,0)--(6.8,0);
  \foreach \i in {-6,...,6}{\draw[muted] (\i,0.12)--(\i,-0.12) node[below=-1pt,font=\scriptsize,text=muted]{$\i$};}
  \draw[accent,line width=3pt] (#1,0)--(#2,0);
  \ifx#3i\fill[accent] (#1-0.25,0)--(#1+0.15,0.2)--(#1+0.15,-0.2)--cycle;\else\extrem{#1}{#3}\fi
  \ifx#4i\fill[accent] (#2+0.25,0)--(#2-0.15,0.2)--(#2-0.15,-0.2)--cycle;\else\extrem{#2}{#4}\fi
  \end{tikzpicture}}

\newcolumntype{C}[1]{>{\centering\arraybackslash}p{#1}}
\newcolumntype{M}[1]{>{\centering\arraybackslash}m{#1}}

\def\fmneed{14}
% =====================================================================
% Capçalera FentMates de les fitxes: un color per a cada curs.
% Els colors són els de la pàgina inicial (LEVELS a assets/js/materials.js);
% el logo ve de _eines/logo/pdf/ (python3 _eines/logo/genera.py).
%
% Ús, al final del preàmbul de la unitat:
%   % =====================================================================
% Capçalera FentMates de les fitxes: un color per a cada curs.
% Els colors són els de la pàgina inicial (LEVELS a assets/js/materials.js);
% el logo ve de _eines/logo/pdf/ (python3 _eines/logo/genera.py).
%
% Ús, al final del preàmbul de la unitat:
%   % =====================================================================
% Capçalera FentMates de les fitxes: un color per a cada curs.
% Els colors són els de la pàgina inicial (LEVELS a assets/js/materials.js);
% el logo ve de _eines/logo/pdf/ (python3 _eines/logo/genera.py).
%
% Ús, al final del preàmbul de la unitat:
%   \input{../comu/fentmates.tex}
%   \fmcurs{1eso}
% i, al començament del document:
%   \capcalera{Matemàtiques}{1r ESO}{Fitxa d'exercicis · Unitat 2}{Divisibilitat}
% Opcional abans de \input: \def\fmneed{14} (espai mínim abans de cada secció).
% =====================================================================
\makeatletter
\def\fm@logodir{../../logo/pdf}
\providecommand{\fmneed}{10}
\def\fm@cursnom{}

% clau del curs, color i color del text damunt del color
\newcommand{\fm@curs}[3]{\@namedef{fm@color@#1}{#2}\@namedef{fm@text@#1}{#3}}
\fm@curs{1eso}{E6B84A}{navy}
\fm@curs{2eso}{3B82F6}{white}
\fm@curs{3eso}{FF7A50}{white}
\fm@curs{4eso-a}{A78BFA}{white}
\fm@curs{4eso-b}{6D28D9}{white}
\fm@curs{1batx-mi}{9B1C1C}{white}
\fm@curs{1batx-ccss}{EF6B6B}{white}
\fm@curs{2batx-mii}{1E6B47}{white}
\fm@curs{2batx-ccss}{2BB3A3}{white}

\newcommand{\fmcurs}[1]{%
  \@ifundefined{fm@color@#1}{\PackageError{fentmates}{Curs desconegut: #1}{}}{}%
  \edef\fm@tmp{\noexpand\definecolor{curs}{HTML}{\@nameuse{fm@color@#1}}%
               \noexpand\colorlet{curstext}{\@nameuse{fm@text@#1}}}\fm@tmp
  \colorlet{cursfosc}{curs!72!black}%
  \def\fm@logo{\fm@logodir/vertical-#1.pdf}%
  \def\fm@marca{\fm@logodir/marca-#1.pdf}}

% ---------- capçalera de la primera pàgina ----------
% #1 matèria (petit), #2 curs, #3 tipus de document i unitat, #4 títol
\newcommand{\capcalera}[4]{%
  \gdef\fm@cursnom{#2}%
  \par\noindent
  \begin{tcolorbox}[colback=curs,colframe=curs,boxrule=0pt,arc=2.2mm,width=3.5cm,height=2.15cm,
      valign=center,left=3mm,right=2mm,top=1mm,bottom=1mm,nobeforeafter,box align=center]
    \color{curstext}{\footnotesize #1}\par\vspace{1pt}{\fontsize{19}{21}\selectfont\bfseries #2\par}%
  \end{tcolorbox}\hspace{5mm}%
  \begin{minipage}[c]{\dimexpr\linewidth-3.5cm-5mm-2.9cm\relax}
    {\color{cursfosc}\small\bfseries\MakeUppercase{#3}}\par\vspace{2pt}
    {\color{navy}\fontsize{20}{23}\selectfont\bfseries #4\par}
  \end{minipage}\hfill
  \begin{minipage}[c]{2.6cm}\raggedleft\includegraphics[width=2.6cm]{\fm@logo}\end{minipage}\par
  \vspace{2pt}}

% ---------- barres de secció del color del curs ----------
\renewcommand{\seccio}[2]{\par\needspace{\fmneed\baselineskip}\bigskip
  \begin{tcolorbox}[colback=curs,colframe=curs,coltext=curstext,arc=2mm,boxrule=0pt,left=3mm,right=3mm,top=1.2mm,bottom=1.2mm]
  {\large\bfseries #1\quad #2}
  \end{tcolorbox}\vspace{-2pt}}

% ---------- capçalera i peu de les pàgines ----------
% primera pàgina: logos institucionals al peu
\fancypagestyle{primera}{\fancyhf{}\renewcommand{\headrulewidth}{0pt}
  \fancyfoot[L]{\raisebox{-0.32\height}{\includegraphics[height=0.78cm]{logo-gva.png}}\hspace{5mm}%
    \raisebox{-0.32\height}{\includegraphics[height=0.78cm]{logo-vermellar.png}}}
  \fancyfoot[R]{\scriptsize\color{muted}\href{https://creativecommons.org/licenses/by-nc-sa/4.0/deed.ca}{\ccbyncsa}\ \ CC BY-NC-SA 4.0 · fentmates.org\qquad\small\thepage}}
% la resta: el símbol a l'esquerra i el curs a la dreta, amb el títol al centre com abans
\fancyhead[L]{\raisebox{-0.3\height}{\includegraphics[height=0.62cm]{\fm@marca}}}
\fancyhead[R]{\small\bfseries\color{cursfosc}\fm@cursnom}
\renewcommand{\headrule}{\color{curs!45}\hrule height \headrulewidth}
\makeatother

%   \fmcurs{1eso}
% i, al començament del document:
%   \capcalera{Matemàtiques}{1r ESO}{Fitxa d'exercicis · Unitat 2}{Divisibilitat}
% Opcional abans de \input: \def\fmneed{14} (espai mínim abans de cada secció).
% =====================================================================
\makeatletter
\def\fm@logodir{../../logo/pdf}
\providecommand{\fmneed}{10}
\def\fm@cursnom{}

% clau del curs, color i color del text damunt del color
\newcommand{\fm@curs}[3]{\@namedef{fm@color@#1}{#2}\@namedef{fm@text@#1}{#3}}
\fm@curs{1eso}{E6B84A}{navy}
\fm@curs{2eso}{3B82F6}{white}
\fm@curs{3eso}{FF7A50}{white}
\fm@curs{4eso-a}{A78BFA}{white}
\fm@curs{4eso-b}{6D28D9}{white}
\fm@curs{1batx-mi}{9B1C1C}{white}
\fm@curs{1batx-ccss}{EF6B6B}{white}
\fm@curs{2batx-mii}{1E6B47}{white}
\fm@curs{2batx-ccss}{2BB3A3}{white}

\newcommand{\fmcurs}[1]{%
  \@ifundefined{fm@color@#1}{\PackageError{fentmates}{Curs desconegut: #1}{}}{}%
  \edef\fm@tmp{\noexpand\definecolor{curs}{HTML}{\@nameuse{fm@color@#1}}%
               \noexpand\colorlet{curstext}{\@nameuse{fm@text@#1}}}\fm@tmp
  \colorlet{cursfosc}{curs!72!black}%
  \def\fm@logo{\fm@logodir/vertical-#1.pdf}%
  \def\fm@marca{\fm@logodir/marca-#1.pdf}}

% ---------- capçalera de la primera pàgina ----------
% #1 matèria (petit), #2 curs, #3 tipus de document i unitat, #4 títol
\newcommand{\capcalera}[4]{%
  \gdef\fm@cursnom{#2}%
  \par\noindent
  \begin{tcolorbox}[colback=curs,colframe=curs,boxrule=0pt,arc=2.2mm,width=3.5cm,height=2.15cm,
      valign=center,left=3mm,right=2mm,top=1mm,bottom=1mm,nobeforeafter,box align=center]
    \color{curstext}{\footnotesize #1}\par\vspace{1pt}{\fontsize{19}{21}\selectfont\bfseries #2\par}%
  \end{tcolorbox}\hspace{5mm}%
  \begin{minipage}[c]{\dimexpr\linewidth-3.5cm-5mm-2.9cm\relax}
    {\color{cursfosc}\small\bfseries\MakeUppercase{#3}}\par\vspace{2pt}
    {\color{navy}\fontsize{20}{23}\selectfont\bfseries #4\par}
  \end{minipage}\hfill
  \begin{minipage}[c]{2.6cm}\raggedleft\includegraphics[width=2.6cm]{\fm@logo}\end{minipage}\par
  \vspace{2pt}}

% ---------- barres de secció del color del curs ----------
\renewcommand{\seccio}[2]{\par\needspace{\fmneed\baselineskip}\bigskip
  \begin{tcolorbox}[colback=curs,colframe=curs,coltext=curstext,arc=2mm,boxrule=0pt,left=3mm,right=3mm,top=1.2mm,bottom=1.2mm]
  {\large\bfseries #1\quad #2}
  \end{tcolorbox}\vspace{-2pt}}

% ---------- capçalera i peu de les pàgines ----------
% primera pàgina: logos institucionals al peu
\fancypagestyle{primera}{\fancyhf{}\renewcommand{\headrulewidth}{0pt}
  \fancyfoot[L]{\raisebox{-0.32\height}{\includegraphics[height=0.78cm]{logo-gva.png}}\hspace{5mm}%
    \raisebox{-0.32\height}{\includegraphics[height=0.78cm]{logo-vermellar.png}}}
  \fancyfoot[R]{\scriptsize\color{muted}\href{https://creativecommons.org/licenses/by-nc-sa/4.0/deed.ca}{\ccbyncsa}\ \ CC BY-NC-SA 4.0 · fentmates.org\qquad\small\thepage}}
% la resta: el símbol a l'esquerra i el curs a la dreta, amb el títol al centre com abans
\fancyhead[L]{\raisebox{-0.3\height}{\includegraphics[height=0.62cm]{\fm@marca}}}
\fancyhead[R]{\small\bfseries\color{cursfosc}\fm@cursnom}
\renewcommand{\headrule}{\color{curs!45}\hrule height \headrulewidth}
\makeatother

%   \fmcurs{1eso}
% i, al començament del document:
%   \capcalera{Matemàtiques}{1r ESO}{Fitxa d'exercicis · Unitat 2}{Divisibilitat}
% Opcional abans de \input: \def\fmneed{14} (espai mínim abans de cada secció).
% =====================================================================
\makeatletter
\def\fm@logodir{../../logo/pdf}
\providecommand{\fmneed}{10}
\def\fm@cursnom{}

% clau del curs, color i color del text damunt del color
\newcommand{\fm@curs}[3]{\@namedef{fm@color@#1}{#2}\@namedef{fm@text@#1}{#3}}
\fm@curs{1eso}{E6B84A}{navy}
\fm@curs{2eso}{3B82F6}{white}
\fm@curs{3eso}{FF7A50}{white}
\fm@curs{4eso-a}{A78BFA}{white}
\fm@curs{4eso-b}{6D28D9}{white}
\fm@curs{1batx-mi}{9B1C1C}{white}
\fm@curs{1batx-ccss}{EF6B6B}{white}
\fm@curs{2batx-mii}{1E6B47}{white}
\fm@curs{2batx-ccss}{2BB3A3}{white}

\newcommand{\fmcurs}[1]{%
  \@ifundefined{fm@color@#1}{\PackageError{fentmates}{Curs desconegut: #1}{}}{}%
  \edef\fm@tmp{\noexpand\definecolor{curs}{HTML}{\@nameuse{fm@color@#1}}%
               \noexpand\colorlet{curstext}{\@nameuse{fm@text@#1}}}\fm@tmp
  \colorlet{cursfosc}{curs!72!black}%
  \def\fm@logo{\fm@logodir/vertical-#1.pdf}%
  \def\fm@marca{\fm@logodir/marca-#1.pdf}}

% ---------- capçalera de la primera pàgina ----------
% #1 matèria (petit), #2 curs, #3 tipus de document i unitat, #4 títol
\newcommand{\capcalera}[4]{%
  \gdef\fm@cursnom{#2}%
  \par\noindent
  \begin{tcolorbox}[colback=curs,colframe=curs,boxrule=0pt,arc=2.2mm,width=3.5cm,height=2.15cm,
      valign=center,left=3mm,right=2mm,top=1mm,bottom=1mm,nobeforeafter,box align=center]
    \color{curstext}{\footnotesize #1}\par\vspace{1pt}{\fontsize{19}{21}\selectfont\bfseries #2\par}%
  \end{tcolorbox}\hspace{5mm}%
  \begin{minipage}[c]{\dimexpr\linewidth-3.5cm-5mm-2.9cm\relax}
    {\color{cursfosc}\small\bfseries\MakeUppercase{#3}}\par\vspace{2pt}
    {\color{navy}\fontsize{20}{23}\selectfont\bfseries #4\par}
  \end{minipage}\hfill
  \begin{minipage}[c]{2.6cm}\raggedleft\includegraphics[width=2.6cm]{\fm@logo}\end{minipage}\par
  \vspace{2pt}}

% ---------- barres de secció del color del curs ----------
\renewcommand{\seccio}[2]{\par\needspace{\fmneed\baselineskip}\bigskip
  \begin{tcolorbox}[colback=curs,colframe=curs,coltext=curstext,arc=2mm,boxrule=0pt,left=3mm,right=3mm,top=1.2mm,bottom=1.2mm]
  {\large\bfseries #1\quad #2}
  \end{tcolorbox}\vspace{-2pt}}

% ---------- capçalera i peu de les pàgines ----------
% primera pàgina: logos institucionals al peu
\fancypagestyle{primera}{\fancyhf{}\renewcommand{\headrulewidth}{0pt}
  \fancyfoot[L]{\raisebox{-0.32\height}{\includegraphics[height=0.78cm]{logo-gva.png}}\hspace{5mm}%
    \raisebox{-0.32\height}{\includegraphics[height=0.78cm]{logo-vermellar.png}}}
  \fancyfoot[R]{\scriptsize\color{muted}\href{https://creativecommons.org/licenses/by-nc-sa/4.0/deed.ca}{\ccbyncsa}\ \ CC BY-NC-SA 4.0 · fentmates.org\qquad\small\thepage}}
% la resta: el símbol a l'esquerra i el curs a la dreta, amb el títol al centre com abans
\fancyhead[L]{\raisebox{-0.3\height}{\includegraphics[height=0.62cm]{\fm@marca}}}
\fancyhead[R]{\small\bfseries\color{cursfosc}\fm@cursnom}
\renewcommand{\headrule}{\color{curs!45}\hrule height \headrulewidth}
\makeatother

\fmcurs{3eso}

\begin{document}
\thispagestyle{primera}

% ---------- encapçalament ----------
\capcalera{Matemàtiques}{3r ESO}{Fitxa d'exercicis · Unitat 1}{Nombres racionals}
\vspace{10pt}
Nom i cognoms: \rule{8.4cm}{0.4pt}\hfill Grup: \rule{1.6cm}{0.4pt}\hfill Data: \rule{2.2cm}{0.4pt}
\vspace{8pt}

\begin{tcolorbox}[colback=soft,colframe=soft,arc=2mm,left=3mm,right=3mm,top=1.5mm,bottom=1.5mm]
\small Exercicis graduats per nivells de dificultat:\quad
\estrelles{1}\ bàsic \qquad \estrelles{2}\ mitjà \qquad \estrelles{3}\ avançat.\\[2pt]
Resol-los al quadern mostrant tots els passos. Simplifica sempre els resultats. Les rectes i les taules es poden completar directament en la fitxa.
\end{tcolorbox}

% =====================================================
\seccio{1}{Conjunts numèrics: naturals, enters i racionals}

\ex{1} Marca amb una creu els conjunts als quals pertany cada nombre:\par\smallskip
{\renewcommand{\arraystretch}{2.1}\centering
\begin{tabular}{|M{2.6cm}|M{1.3cm}|M{1.3cm}|M{1.3cm}|}
\hline
\rowcolor{soft}\textbf{Nombre} & $\mathbb{N}$ & $\mathbb{Z}$ & $\mathbb{Q}$ \\ \hline
$15$ & & & \\ \hline
$-8$ & & & \\ \hline
$\dfrac34$ & & & \\ \hline
$0$ & & & \\ \hline
$-\dfrac{12}{3}$ & & & \\ \hline
$0,6$ & & & \\ \hline
$1,\overline{3}$ & & & \\ \hline
$\dfrac{20}{4}$ & & & \\ \hline
$-2,5$ & & & \\ \hline
\end{tabular}\par}

\ex{2} Indica el conjunt numèric més xicotet al qual pertany cada nombre:
\begin{tasks}(4)
  \task $-\dfrac{36}{9}$
  \task $\dfrac{42}{6}$
  \task $-\dfrac72$
  \task $0,125$
  \task $-(-5)$
  \task $\dfrac{18}{-3}$
  \task $2,\overline{7}$
  \task $(-2)^{2}$
\end{tasks}

\ex{2} Indica si les afirmacions següents són vertaderes o falses. Justifica-ho.
\begin{tasks}(2)
  \task Tot nombre enter és racional.
  \task Tot nombre racional és enter.
  \task El nombre $-3$ és natural.
  \task La resta de dos naturals sempre és un natural.
  \task Entre $\frac13$ i $\frac12$ hi ha algun nombre racional.
  \task La divisió de dos enters (amb divisor no nul) sempre és un racional.
\end{tasks}

\ex{3} Respon:
\begin{tasks}(1)
  \task Escriu tres nombres racionals compresos entre $\frac14$ i $\frac13$.
  \task Quants nombres enters hi ha entre $-3,5$ i $2,1$? Escriu-los.
  \task Escriu un nombre que siga enter però no natural, i un altre que siga racional però no enter.
\end{tasks}

% =====================================================
\seccio{2}{Nombres enters: valor absolut, oposat i ordre}

\ex{1} Calcula:
\begin{tasks}(4)
  \task $|-12|$
  \task $|9|$
  \task $\text{op}(-7)$
  \task $\text{op}(15)$
  \task $|\text{op}(-4)|$
  \task $-|-6|$
  \task $\text{op}\big(|-3|\big)$
  \task $|0|$
\end{tasks}

\ex{1} Respon:
\begin{tasks}(1)
  \task Ordena de menor a major: \ $-3,\ \ 5,\ \ -8,\ \ 0,\ \ 2,\ \ -1,\ \ -12$.
  \task Ordena de major a menor: \ $7,\ \ -7,\ \ 4,\ \ -10,\ \ -2,\ \ 1$.
  \task Representa a la recta els nombres $-5$, $3$ i $-2$ i també els seus oposats.\par\smallskip
  \rectabuida
\end{tasks}

\ex{2} Completa amb el signe $<$ o $>$:
\begin{tasks}(3)
  \task $-5\ \square\ -2$
  \task $|-5|\ \square\ |-2|$
  \task $-10\ \square\ 0$
  \task $\text{op}(3)\ \square\ \text{op}(-3)$
  \task $-(-4)\ \square\ -4$
  \task $|-7|\ \square\ 6$
\end{tasks}

\ex{2} Escriu tots els nombres enters $x$ que compleixen:
\begin{tasks}(2)
  \task $|x|=8$
  \task $|x|<4$
  \task $|x|\le 2$ \ i \ $x$ negatiu
  \task $-4<x\le 1$
\end{tasks}

\ex{3} Les temperatures mínimes d'una setmana a Morella han sigut: dilluns $-3\,^\circ$C, dimarts $2\,^\circ$C, dimecres $-5\,^\circ$C, dijous $0\,^\circ$C i divendres $-1\,^\circ$C.
\begin{tasks}(1)
  \task Ordena els dies del més fred al més calorós.
  \task Quina diferència hi ha entre la temperatura més alta i la més baixa?
  \task Quin dia la temperatura va tindre el major valor absolut?
\end{tasks}

% =====================================================
\seccio{3}{Suma i resta de nombres enters}

\ex{1} Calcula:
\begin{tasks}(3)
  \task $(+7)+(-4)$
  \task $(-6)+(-9)$
  \task $(-12)+(+5)$
  \task $(+3)-(-8)$
  \task $(-4)-(+6)$
  \task $(-2)-(-2)$
\end{tasks}

\ex{1} Lleva els parèntesis i calcula:
\begin{tasks}(2)
  \task $8-(-3)+(-5)$
  \task $-4+(-7)-(-10)$
  \task $-(-6)-(+9)+2$
  \task $15+(-20)-(-5)-(+3)$
\end{tasks}

\ex{2} Calcula agrupant els positius i els negatius:
\begin{tasks}(2)
  \task $-3+8-12+5-1$
  \task $10-15+7-9+11$
  \task $-20+13-6+18-5$
  \task $4-7-8+2-3+10$
\end{tasks}

\ex{2} Calcula, començant pels parèntesis de dins:
\begin{tasks}(2)
  \task $7-(3-9)+(-4+1)$
  \task $-(5-8)-(2+6)+4$
  \task $12-\left[5-(3-7)\right]$
  \task $-6-\left[-(2-9)+4\right]-1$
  \task $\left[8-(-3+5)\right]-\left[-4-(6-10)\right]$
\end{tasks}

\ex{3} Resol aquests problemes:
\begin{tasks}(1)
  \task Un submarí és a $120$ m de profunditat. Puja $45$ m, després baixa $30$ m i finalment puja $70$ m. A quina profunditat es troba ara?
  \task Arquimedes va nàixer l'any $287$ aC i va morir l'any $212$ aC. Quants anys va viure?
  \task A les 6 h el termòmetre marca $-4\,^\circ$C. A les 14 h la temperatura ha pujat $11$ graus i a les 22 h ha baixat $9$ graus respecte de les 14 h. Quina temperatura fa a les 22 h?
\end{tasks}

% =====================================================
\seccio{4}{Multiplicació i divisió de nombres enters}

\ex{1} Calcula:
\begin{tasks}(4)
  \task $(-5)\cdot(+6)$
  \task $(-4)\cdot(-9)$
  \task $(+7)\cdot(-8)$
  \task $(-1)\cdot(-1)$
  \task $(-48):(-6)$
  \task $(+54):(-9)$
  \task $(-63):(+7)$
  \task $0:(-5)$
\end{tasks}

\ex{1} Calcula aplicant la regla dels signes:
\begin{tasks}(2)
  \task $(-2)\cdot3\cdot(-5)$
  \task $(-1)\cdot(-2)\cdot(-3)\cdot(-4)$
  \task $(-2)\cdot(-2)\cdot(-2)\cdot5$
  \task $(-3)\cdot(-10)\cdot0\cdot(-7)$
\end{tasks}

\ex{2} Completa el nombre que falta:
\begin{tasks}(4)
  \task $(-6)\cdot\square=42$
  \task $\square:(-4)=9$
  \task $56:\square=-8$
  \task $(-3)\cdot\square\cdot(-2)=-30$
\end{tasks}

\ex{2} Calcula d'esquerra a dreta:
\begin{tasks}(2)
  \task $(-72):(-6):(-3)$
  \task $(-60):(-5)\cdot(-2)$
  \task $100:(-4):(-5)$
  \task $(-2)\cdot(-18):(-4)$
\end{tasks}

\ex{3} Respon:
\begin{tasks}(1)
  \task El producte de dos nombres enters és $-36$ i la seua suma és $0$. Quins nombres són?
  \task Escriu tres parelles de nombres enters el producte dels quals siga $-12$.
  \task Sense calcular cap potència, quin és el resultat de $(-1)^{2025}\cdot(-1)^{2026}$?
\end{tasks}

% =====================================================
\seccio{5}{Operacions combinades amb nombres enters}

\ex{1} Calcula respectant la jerarquia de les operacions:
\begin{tasks}(2)
  \task $5+3\cdot(-2)$
  \task $-8-12:(-4)$
  \task $4\cdot(-3)-2\cdot(-5)$
  \task $20:(-5)+(-3)\cdot(-6)$
\end{tasks}

\ex{2} Calcula:
\begin{tasks}(2)
  \task $6-2\cdot(5-9)$
  \task $(7-12):(-5)+3$
  \task $-3\cdot(4-6)^{2}$
  \task $18:(-3)\cdot2-(-4)$
  \task $(-2)^{3}+3^{2}-(-1)^{5}$
  \task $-2^{4}+(-2)^{4}$
\end{tasks}

\ex{2} Calcula:
\begin{tasks}(2)
  \task $15-\left[4-2\cdot(3-8)\right]$
  \task $-3\cdot\left[(-2)^{2}-7\right]+10:(-2)$
  \task $\left[(-6)\cdot3+20\right]:(-2)-(-5)$
  \task $2\cdot\left\{5-\left[3-(4-9)\right]\right\}$
\end{tasks}

\ex{3} Calcula:
\begin{tasks}(1)
  \task $(-3)^{2}-4\cdot\left[6-(-2)^{3}:(-4)\right]+(-10):5$
  \task $\left\{-2\cdot\left[3-(-1)^{4}\cdot5\right]+(-3)^{2}\right\}:(-13)$
  \task $25:(-5)^{2}-\left[(-4)\cdot(-3)-2^{4}\right]\cdot(-1)^{3}$
  \task Col·loca un parèntesi en $\ 3\cdot4-6:2+1\ $ perquè el resultat siga $-2$. I perquè siga $4$?
\end{tasks}

% =====================================================
\seccio{6}{Fraccions i nombres racionals}

\ex{1} Simplifica fins a obtindre la fracció irreductible:
\begin{tasks}(6)
  \task $\dfrac{12}{18}$
  \task $\dfrac{45}{60}$
  \task $\dfrac{28}{35}$
  \task $-\dfrac{40}{56}$
  \task $\dfrac{66}{121}$
  \task $\dfrac{84}{36}$
\end{tasks}

\ex{1} Comprova si són equivalents:
\begin{tasks}(4)
  \task $\dfrac35$\ i\ $\dfrac{12}{20}$
  \task $\dfrac46$\ i\ $\dfrac{10}{15}$
  \task $\dfrac78$\ i\ $\dfrac{21}{32}$
  \task $-\dfrac29$\ i\ $\dfrac{6}{-27}$
\end{tasks}

\ex{2} Calcula el valor de $x$ perquè les fraccions siguen equivalents:
\begin{tasks}(4)
  \task $\dfrac23=\dfrac{x}{15}$
  \task $\dfrac5x=\dfrac{20}{36}$
  \task $\dfrac{x}{7}=\dfrac{18}{42}$
  \task $\dfrac{12}{16}=\dfrac{9}{x}$
\end{tasks}

\ex{2} Redueix a comú denominador i ordena de menor a major:
\begin{tasks}(1)
  \task $\dfrac34,\ \ \dfrac56,\ \ \dfrac23,\ \ \dfrac{7}{12}$
  \task $-\dfrac12,\ \ \dfrac35,\ \ -\dfrac34,\ \ 0,\ \ \dfrac{7}{10}$
\end{tasks}

\ex{3} Respon:
\begin{tasks}(1)
  \task Escriu una fracció de denominador $20$ compresa entre $\frac25$ i $\frac35$.
  \task Quina fracció d'una hora són $25$ minuts? Simplifica-la.
  \task Troba una fracció equivalent a $\frac37$ en què la suma del numerador i el denominador siga $60$.
\end{tasks}

% =====================================================
\seccio{7}{Operacions amb fraccions}

\ex{1} Calcula i simplifica:
\begin{tasks}(3)
  \task $\dfrac37+\dfrac27$
  \task $\dfrac59-\dfrac89$
  \task $\dfrac12+\dfrac13$
  \task $\dfrac34-\dfrac56$
  \task $2-\dfrac35$
  \task $\dfrac16+\dfrac29-\dfrac13$
\end{tasks}

\ex{1} Calcula i simplifica:
\begin{tasks}(3)
  \task $\dfrac23\cdot\dfrac94$
  \task $\left(-\dfrac56\right)\cdot\dfrac{3}{10}$
  \task $4\cdot\left(-\dfrac78\right)$
  \task $\dfrac35:\dfrac{9}{10}$
  \task $\left(-\dfrac47\right):\left(-\dfrac{2}{21}\right)$
  \task $6:\dfrac34$
\end{tasks}

\ex{2} Calcula i simplifica:
\begin{tasks}(2)
  \task $\dfrac{5}{12}+\dfrac38-\dfrac16$
  \task $-\dfrac{2}{15}-\dfrac{3}{10}+\dfrac16$
  \task $3-\dfrac25-\dfrac{7}{10}$
  \task $\dfrac14-\left(\dfrac23-\dfrac56\right)$
\end{tasks}

\ex{2} Calcula simplificant abans de multiplicar:
\begin{tasks}(2)
  \task $\dfrac{14}{15}\cdot\dfrac{25}{28}$
  \task $\left(-\dfrac{18}{35}\right)\cdot\dfrac{21}{12}\cdot\dfrac{10}{9}$
  \task $\dfrac{12}{25}:\dfrac{18}{35}$
  \task $\left(-\dfrac{8}{27}\right):\left(-\dfrac{16}{9}\right)\cdot\dfrac32$
\end{tasks}

\ex{3} Respon:
\begin{tasks}(1)
  \task Quina fracció cal sumar a $\frac38$ per a obtindre $\frac56$?
  \task Per quina fracció cal multiplicar $\frac49$ per a obtindre $\frac23$?
  \task La suma de dos nombres és $\frac{7}{12}$ i un d'ells és $-\frac14$. Quin és l'altre?
  \task Quin és l'invers de l'oposat de $\frac27$?
\end{tasks}

% =====================================================
\seccio{8}{Potències de fraccions}

\ex{1} Calcula:
\begin{tasks}(4)
  \task $\left(\dfrac34\right)^{2}$
  \task $\left(-\dfrac25\right)^{3}$
  \task $\left(-\dfrac13\right)^{4}$
  \task $-\left(\dfrac12\right)^{4}$
  \task $\left(\dfrac57\right)^{0}$
  \task $\left(\dfrac23\right)^{-2}$
  \task $\left(-\dfrac14\right)^{-1}$
  \task $\left(\dfrac32\right)^{-3}$
\end{tasks}

\ex{2} Expressa com una sola potència i, si es pot, calcula:
\begin{tasks}(3)
  \task $\left(\dfrac23\right)^{4}\cdot\left(\dfrac23\right)^{3}$
  \task $\left(\dfrac54\right)^{6}:\left(\dfrac54\right)^{2}$
  \task $\left[\left(-\dfrac12\right)^{3}\right]^{2}$
  \task $\left(\dfrac35\right)^{4}\cdot\left(\dfrac53\right)^{4}$
  \task $\left(\dfrac72\right)^{5}:\left(\dfrac72\right)^{5}$
  \task $\left(\dfrac49\right)^{3}:\left(\dfrac23\right)^{3}$
\end{tasks}

\ex{2} Calcula:
\begin{tasks}(2)
  \task $\left(\dfrac12\right)^{2}+\left(\dfrac12\right)^{3}$
  \task $\left(\dfrac23\right)^{2}-\left(\dfrac13\right)^{2}$
  \task $\left(-\dfrac32\right)^{2}\cdot\left(\dfrac23\right)^{3}$
  \task $\left(-\dfrac12\right)^{3}:\left(\dfrac14\right)^{2}$
\end{tasks}

\ex{3} Simplifica i calcula:
\begin{tasks}(2)
  \task $\dfrac{\left(\frac23\right)^{5}\cdot\left(\frac32\right)^{2}}{\left(\frac23\right)^{2}}$
  \task $\left[\left(\dfrac12\right)^{3}\cdot\left(\dfrac12\right)^{-5}\right]^{2}$
  \task $\left[\left(\dfrac35\right)^{2}\right]^{3}:\left(\dfrac35\right)^{5}$
  \task $\left(-\dfrac23\right)^{4}\cdot\left(-\dfrac32\right)^{3}$
\end{tasks}

% =====================================================
\seccio{9}{Operacions combinades amb fraccions}

\ex{1} Calcula i simplifica:
\begin{tasks}(2)
  \task $\dfrac12+\dfrac13\cdot\dfrac34$
  \task $\dfrac23-\dfrac16:2$
  \task $\dfrac56\cdot\left(\dfrac12+\dfrac13\right)$
  \task $\left(\dfrac34-\dfrac12\right):\dfrac58$
\end{tasks}

\ex{2} Calcula i simplifica:
\begin{tasks}(2)
  \task $\dfrac23+\dfrac12\cdot\left(\dfrac53-1\right)$
  \task $\left(1-\dfrac34\right)\cdot\left(2+\dfrac23\right)$
  \task $\dfrac35-\dfrac25:\left(\dfrac12-\dfrac13\right)$
  \task $\left(\dfrac12\right)^{2}+\dfrac34:\left(\dfrac32-1\right)$
\end{tasks}

\ex{2} Calcula i simplifica:
\begin{tasks}(2)
  \task $\left[\dfrac23-\left(\dfrac12\right)^{2}\right]:\dfrac56$
  \task $1-\left[\dfrac34-\left(\dfrac12-\dfrac13\right)\right]$
  \task $\left(-\dfrac23\right)^{2}:\dfrac43-\dfrac16$
  \task $\dfrac32-\dfrac23\cdot\left[\dfrac14-\left(1-\dfrac12\right)^{2}\cdot3\right]$
\end{tasks}

\ex{3} Calcula i simplifica:
\begin{tasks}(2)
  \task $\left\{\dfrac12-\left[\dfrac23-\left(\dfrac34-\dfrac16\right)\right]\right\}:\left(\dfrac13\right)^{2}$
  \task $\left(1-\dfrac12\right)^{-2}\cdot\left(\dfrac34-1\right)+\left(-\dfrac12\right)^{3}\cdot8$
  \task $\dfrac{\dfrac23-\dfrac12}{1+\dfrac13}$
  \task $\dfrac{2-\dfrac{1}{1+\frac12}}{\dfrac12+\dfrac13}$
\end{tasks}

% =====================================================
\seccio{10}{Problemes amb nombres racionals}

\ex{1} Dels $840$ alumnes d'un institut, $\frac{5}{12}$ viuen al mateix poble on és el centre. Quants alumnes hi viuen? Quants hi venen de fora?

\ex{1} Un ciclista ha recorregut $\frac35$ d'una etapa de $165$ km. Quants quilòmetres ha fet? Quants n'hi falten?

\ex{1} Els $\frac38$ d'una quantitat de diners són $45$ €. Quina és la quantitat total?

\ex{2} En una biblioteca, $\frac25$ dels llibres són novel·les, $\frac14$ són de ciència i la resta, $840$ llibres, són d'altres temes. Quants llibres hi ha en total? Quantes novel·les?

\ex{2} Una mare reparteix $600$ € entre els seus tres fills: el primer en rep $\frac13$, el segon $\frac25$ i el tercer, la resta. Quants diners rep cadascun?

\ex{2} Quants gots de $\frac{3}{16}$ de litre es poden omplir amb $6$ ampolles de $\frac34$ de litre?

\ex{2} D'un dipòsit ple d'aigua es gasta $\frac13$ el dilluns i $\frac34$ del que quedava el dimarts. Quina fracció del dipòsit queda? Si la capacitat és de $1\,200$ L, quants litres queden?

\ex{3} Una persona gasta $\frac25$ del sou en el lloguer, $\frac13$ del que li queda en menjar i estalvia els $600$ € restants. Quin és el seu sou? Quant gasta en cada cosa?

\ex{3} Una aixeta omple un dipòsit en $4$ hores i una altra, en $12$ hores.
\begin{tasks}(1)
  \task Quant de temps tarden a omplir-lo si s'obrin totes dues alhora?
  \task I si, a més, està obert un desguàs que buida el dipòsit ple en $6$ hores?
\end{tasks}

\ex{3} Un llaurador ven el primer dia $\frac38$ de les síndries que té, el segon dia $\frac25$ de les que li queden i el tercer dia, les $60$ que encara li quedaven. Quantes síndries tenia?

\ex{3} Una pilota bota i cada vegada arriba als $\frac23$ de l'altura anterior. Si la deixem caure des de $81$ cm, a quina altura arriba després del tercer bot?

% =====================================================
\seccio{11}{Nombres decimals: classificació}

\ex{1} Classifica cada decimal (exacte, periòdic pur, periòdic mixt o no periòdic) i, si és periòdic, indica'n la part entera, l'anteperíode i el període:
\begin{tasks}(3)
  \task $2,375$
  \task $0,4444\ldots$
  \task $1,2333\ldots$
  \task $3,141\,592\ldots$
  \task $0,121212\ldots$
  \task $5,1060606\ldots$
\end{tasks}

\ex{1} Fes la divisió, escriu l'expressió decimal i classifica-la:
\begin{tasks}(6)
  \task $\dfrac78$
  \task $\dfrac29$
  \task $\dfrac56$
  \task $\dfrac{13}{11}$
  \task $\dfrac{7}{20}$
  \task $\dfrac{11}{15}$
\end{tasks}

\ex{2} Sense fer la divisió, indica quin tipus d'expressió decimal té cada fracció. Justifica-ho.
\begin{tasks}(4)
  \task $\dfrac{3}{40}$
  \task $\dfrac{5}{21}$
  \task $\dfrac{7}{30}$
  \task $\dfrac{9}{12}$
  \task $\dfrac{4}{33}$
  \task $\dfrac{11}{18}$
  \task $\dfrac{14}{35}$
  \task $\dfrac{6}{55}$
\end{tasks}

\ex{3} Respon:
\begin{tasks}(1)
  \task Escriu una fracció de denominador $12$ que tinga una expressió decimal exacta.
  \task Ordena de menor a major: \ $0,\overline{34}$;\ \ $0,3$;\ \ $0,3\overline{4}$;\ \ $0,\overline{3}$.
  \task Les fraccions $\frac13$ i $\frac16$ donen decimals periòdics. Pot ser exacta la seua suma? Comprova-ho.
\end{tasks}

% =====================================================
\seccio{12}{Pas de decimal a fracció}

\ex{1} Troba la fracció generatriu i simplifica-la:
\begin{tasks}(4)
  \task $0,8$
  \task $2,45$
  \task $0,036$
  \task $3,125$
\end{tasks}

\ex{2} Troba la fracció generatriu i simplifica-la:
\begin{tasks}(4)
  \task $0,\overline{7}$
  \task $2,\overline{3}$
  \task $0,\overline{18}$
  \task $1,\overline{36}$
\end{tasks}

\ex{2} Troba la fracció generatriu i simplifica-la:
\begin{tasks}(4)
  \task $0,1\overline{6}$
  \task $1,2\overline{7}$
  \task $0,0\overline{3}$
  \task $2,1\overline{36}$
\end{tasks}

\ex{3} Respon:
\begin{tasks}(1)
  \task Demostra que $0,\overline{9}=1$.
  \task Calcula $0,\overline{3}+0,\overline{6}$ passant abans a fracció.
  \task Quin nombre és més gran, $0,\overline{45}$ o $0,4\overline{5}$? Comprova-ho amb les fraccions generatrius.
  \task Troba la fracció generatriu de $0,\overline{142857}$.
\end{tasks}

% =====================================================
\seccio{13}{Operacions combinades amb fraccions i decimals}

\ex{2} Passa els decimals a fracció i calcula:
\begin{tasks}(2)
  \task $0,5+0,\overline{3}$
  \task $1,2\cdot0,\overline{6}$
  \task $0,\overline{6}-0,1\overline{6}$
  \task $2,\overline{3}:0,\overline{7}$
\end{tasks}

\ex{2} Calcula i dona el resultat en forma de fracció irreductible:
\begin{tasks}(2)
  \task $0,25+1,\overline{3}\cdot0,75$
  \task $\left(0,\overline{3}+0,5\right)^{2}$
  \task $1,5-0,\overline{2}:0,\overline{4}$
  \task $\left(2,4-0,\overline{6}\right)\cdot0,1\overline{6}$
\end{tasks}

\ex{3} Calcula i dona el resultat en forma de fracció irreductible:
\begin{tasks}(2)
  \task $\left(1,\overline{6}-0,5\right):\left(0,\overline{3}\right)^{2}$
  \task $\dfrac{0,\overline{3}+0,25}{1,1\overline{6}}$
  \task $\left(0,\overline{9}-0,1\right)\cdot1,\overline{1}$
  \task $2,0\overline{3}\cdot0,6-0,2$
\end{tasks}

\vfill
\begin{center}\color{muted}\small\faCheckCircle[regular]\ \ Ànim! Si et costa algun exercici, repassa els apunts interactius de la unitat.\end{center}

\end{document}
